\centerline{\bf \Huge Antiparallel parallel club}
\centerline{\bf \Huge }

Пусть фиксирован угол $B A C$. На прямых $A B$ и $A C$ взяты точки $B_1$ и $C_1$ (либо обе на сторонах угла, либо обе на продолжениях сторон). Тогда прямые $B_1 C_1$ и $B C$ называются \textit{антипараллельными} относительно угла $B A C$, если $\angle A B C=\angle A C_1 B_1$. 

\textbf{Очень важное замечание!} Антипараллельность это НЕ взаимотношение двух прямых, как с обычной параллельностью. Тут очень важно наличие фиксированного угла, поэтому в задачах мы всегда уточняем "прямые антипараллельны относительно угла..."\!.

\problem{0.1} В предыдущих обозначениях, $B C$ антипараллельна $B_1 C_1$ относительно угла $BAC$ тогда и только тогда, когда $B, C, B_1, C_1$ лежат на одной окружности.

\problem{0.2} Точки $B_1$ и $C_1$ могут обе совпасть с точкой $A$ : тогда прямая $B_1 C_1$ станет касательной к описанной окружности треугольника $A B C$. Таким образом, касательная к описанной окружности треугольника --- \textit{антипараллель} к соответствующей стороне. 

\problem{1}
Прямые, антипараллельные одной и той же прямой относительно некоторого угла, параллельны.

\problem{2}
Радиус описанной окружности, проведенный из вершины треугольника, перпендикулярен всем прямым, антипараллельным противоположной стороне.

\problem{3}
В остроугольном треугольнике $A B C$ проведены высоты $A A_1$ и $C C_1$. Из точек $A_1$ и $C_1$ опустили перпендикуляры на прямые $A B$ и $B C$. Докажите, что прямая, проходящая через основания этих перпендикуляров, параллельна $A C$.

\problem{4}
\textbf{Теорема Брахмагупты.}
Четырёхугольник, диагонали которого взаимно перпендикулярны, вписан в окружность. Докажите, что продолжение перпендикуляра из точки пересечения диагоналей к одной из сторон делит противоположную сторону пополам.

\problem{5}
\textbf{Лемма Фусса.}
Окружности пересекаются в точках $M$ и $K$. Через $M$ и $K$ проведены прямые $A B$ и $C D$ соответственно, пересекающие первую окружность в точках $A$ и $C$, вторую в точках $B$ и $D$. Докажите, что $A C \| B D$.

\problem{6}
Дан вписанный четырёхугольник. Для каждой вершины рассмотрим её проекцию на диагональ, не содержащую эту вершину. Докажите, что четыре полученные точки лежат на одной окружности.

